% Generated by scripts/gen_blueprint.py from blueprint/nodes/10-rank-one-configurations.toml. Do not edit.

\chapter{Rank-one configurations near a point of a circle}
\label{chap:10-rank-one-configurations}

Results of the library's work on Diaz's conjecture, about the hypothesis that rank-one statements need. The six exponentials theorem (Chapter~\ref{chap:05-six-exponentials}) and Roy's strong six exponentials theorem \cite{Roy92} take as input a \emph{configuration}: $x_1,x_2$ linearly independent over $\Qbar$ and $y_1,y_2,y_3$ linearly independent over $\Qbar$, with all six products $x_iy_j$ in a given $\Qbar$-vector space. Near a point $u$ of a circle of algebraic radius, that is $u\neq0$ with $u\bar u$ algebraic, Theorem~\ref{thm:generic_circle_point_no_two_by_three_configuration} shows that generic data carry no configuration, whatever numbers are added, and Theorem~\ref{thm:circle_point_extension_two_by_three_configuration_iff} classifies the four-dimensional extensions $\Qbar+\Qbar u+\Qbar\bar u+\Qbar z$ that carry one. Neither statement involves $e^{u}$, so both keep their content if Diaz's conjecture holds. Neither was found in the sources read; the printed special cases are named in each source line. A third section adds two families of five-dimensional spaces $\Qbar+\Qbar u+\Qbar\bar u+\Qbar z+\Qbar\bar z$ that carry a configuration invisible to those extensions (Theorem~\ref{thm:circle_point_conjugate_pair_configuration_invisible_to_four_dimensional_extensions}); that was not found either, but it is short, and its consequences at a candidate are printed (Corollary~\ref{cor:candidate_div_sq_sub_not_mem_logAlgTilde}). The last section shows that in the power hulls $\Qbar u^{-k}+\Qbar u^{-1}+\Qbar+\Qbar u+\Qbar u^{k}$ a configuration exists only for $k=2$ and $k=3$ (Theorem~\ref{thm:power_hull_strong_six_exp_configuration_iff}), so that the strong six exponentials route at a candidate stops at $u^{3}$.

\section{Generic data}

% Prove2Me: DiazModulus.generic_circle_point_no_two_by_three_configuration -- https://prove2.me/theorems/2f818ada-bb16-408c-9e9d-66186a87e8fb
\begin{theorem}[No configuration on generic data]
\label{thm:generic_circle_point_no_two_by_three_configuration}
\lean{Diaz.generic_circle_point_no_two_by_three_configuration}
\leanok
Let $u\neq0$ with $u\bar u$ algebraic, and let $w_1,\dots,w_m\in\C$ be such that $u,w_1,\dots,w_m$ are algebraically independent over $\Qbar$. Then there are no $x_1,x_2$, linearly independent over $\Qbar$, and $y_1,y_2,y_3$, linearly independent over $\Qbar$, with all six products $x_iy_j$ in $\Qbar+\Qbar u+\Qbar\bar u+\sum_j\Qbar w_j$.

{\small\emph{Source:} \textbf{Not found in the sources read, and not routine} (Chapter~\ref{chap:14-not-found-in-the-sources}). The case $m=0$ is \cite[Theorem 3.4]{Roy95} and \cite[Lemma 12.16 and Exercise 12.10]{Wal00}; see also \cite[\S3.2]{Roy95} and \cite[p.~186]{Fis01}. The general statement was not found in the sources read.}
\end{theorem}

\begin{proof}
\leanok
Multiply by $u$. Since $\bar u=\rho/u$ with $\rho$ algebraic, every element of the space times $u$ is the value at $(u,w)$ of a polynomial over $\Qbar$ of total degree at most two that is constant once its first variable $X_0$ is set to $0$, and evaluation at $(u,w)$ is injective. A configuration therefore gives a rank-one $2\times3$ matrix of such polynomials, which factors as $(h,g)\otimes(c_0,c_1,c_2)$ with both factors free. Total degrees add, and $X_0\mapsto0$ is a ring map; each case puts two free vectors in a line or three in a plane.
\end{proof}

\section{Four-dimensional extensions}

% Prove2Me: DiazModulus.two_by_three_configuration_forces_progression -- https://prove2.me/theorems/46109c0d-3af4-46c3-900b-441e398cde5e
\begin{lemma}[A configuration in dimension four is a progression]
\label{lem:two_by_three_configuration_forces_progression}
\lean{Diaz.two_by_three_configuration_forces_progression}
\leanok
Let $V\subseteq\C$ be a $\Qbar$-vector subspace of dimension at most $4$, and let $x_1,x_2$ and $y_1,y_2,y_3$ be linearly independent over $\Qbar$ with all six products $x_iy_j$ in $V$. Then
\[ V=\Qbar\,b+\Qbar\,bh+\Qbar\,bh^{2}+\Qbar\,bh^{3} \]
for some $b\neq0$ and $h\notin\Qbar$.

{\small\emph{Source:} The converse direction, that such a progression carries a configuration, is used in \cite[Lemma 6.1]{Fis01} and \cite[Théorème 7(2)]{Dia07}. This direction was not found in the sources read.}
\end{lemma}

\begin{proof}
\leanok
Put $h=x_2/x_1$, which is not in $\Qbar$ and so is transcendental, and $B=x_1(\Qbar y_1+\Qbar y_2+\Qbar y_3)$. Both $B$ and $hB$ lie in $V$, so $W=B\cap hB$ has dimension at least $2$. The spaces $W$ and $h^{-1}W$ lie in $B$ and have dimension at least $2$, so they meet in some $c=hb\neq0$ with $b\in B$. Then $b,hb,h^{2}b\in B$ and $h^{3}b\in hB$: four free vectors in $V$.
\end{proof}

% Prove2Me: DiazModulus.cubic_product_eq_square_normal_form -- https://prove2.me/theorems/bbd5da86-854b-4c2d-b59d-2fe6ec883487
\begin{lemma}[A normal form for $P_1P_2=P_0^2$]
\label{lem:cubic_product_eq_square_normal_form}
\lean{Diaz.cubic_product_eq_square_normal_form}
\leanok
Let $K$ be a field and $P_0,P_1,P_2\in K[X]$ with $\deg P_1,\deg P_2\le3$, $P_1P_2=P_0^{2}$, and $P_0,P_1$ linearly independent over $K$. Then there are $Q_0,Q_1\in K[X]$ of degree at most $1$, linearly independent over $K$, and $a,b\in K$ such that, with $g=aQ_0+bQ_1$,
\[ P_0=g\,Q_0Q_1,\qquad P_1=g\,Q_1^{2},\qquad P_2=g\,Q_0^{2}. \]

{\small\emph{Source:} Elementary.}
\end{lemma}

\begin{proof}
\leanok
Divide out a common factor: $P_0=dQ_0$ and $P_1=dQ_1$ with $Q_0,Q_1$ coprime. Then $Q_1P_2=dQ_0^{2}$, so $Q_1$ divides $d$, say $d=gQ_1$. Degrees add, so $\deg Q_i\le1$. A relation between $Q_0$ and $Q_1$ would be one between $P_0$ and $P_1$, so they are independent; one of them has degree $1$, hence $\deg g\le1$, and by Cramer's rule $g=aQ_0+bQ_1$.
\end{proof}

% Prove2Me: DiazModulus.circle_hull_progression_contains_square_or_reciprocal -- https://prove2.me/theorems/1706d5e7-e333-42e4-8cad-2b8657fafca8
\begin{lemma}[A progression through $1$, $u$ and $\bar u$]
\label{lem:circle_hull_progression_contains_square_or_reciprocal}
\lean{Diaz.circle_hull_progression_contains_square_or_reciprocal}
\leanok
Let $u\in\C\setminus\Qbar$ with $u\bar u$ algebraic, let $b\neq0$ and $h\notin\Qbar$, and suppose that $V=\Qbar\,b+\Qbar\,bh+\Qbar\,bh^{2}+\Qbar\,bh^{3}$ contains $1$, $u$ and $\bar u$. Then $V$ contains $u^{2}$, or $\bar u^{2}$, or $1/(u-a)$ for some $a\in\Qbar$, $a\neq0$.

{\small\emph{Source:} A step of Theorem~\ref{thm:circle_point_extension_two_by_three_configuration_iff}.}
\end{lemma}

\begin{proof}
\leanok
\uses{lem:cubic_product_eq_square_normal_form}
Write $1=bP_0(h)$, $u=bP_1(h)$ and $u^{-1}=\bar u/(u\bar u)=bP_2(h)$ with $\deg P_i\le3$. As $h$ is transcendental, $P_1P_2=P_0^{2}$, and $P_0,P_1$ are independent because $u\notin\Qbar$. Lemma~\ref{lem:cubic_product_eq_square_normal_form} gives $g=aQ_0+cQ_1$ and $u=Q_1(h)/Q_0(h)$. If $c=0$, then $u^{2}=b\,aQ_1(h)^{3}$; if $a=0$, then $\bar u^{2}=b\,\rho^{2}c\,Q_0(h)^{3}$ with $\rho=u\bar u$; otherwise $1/(u+a/c)=b\,c\,Q_0(h)^{2}Q_1(h)$.
\end{proof}

% Prove2Me: DiazModulus.circle_point_extension_carries_two_by_three_configuration -- https://prove2.me/theorems/dd39ce78-249e-4099-a7ff-76acfc50ff77
\begin{proposition}[Three spaces that carry a configuration]
\label{prop:circle_point_extension_carries_two_by_three_configuration}
\lean{Diaz.circle_point_extension_carries_two_by_three_configuration}
\leanok
Let $u\in\C\setminus\Qbar$ with $u\bar u$ algebraic, and let $w$ be $u^{2}$, $\bar u^{2}$, or $1/(u-a)$ with $a\in\Qbar$, $a\neq0$. Then there are $x_1,x_2$, linearly independent over $\Qbar$, and $y_1,y_2,y_3$, linearly independent over $\Qbar$, with all six products $x_iy_j$ in $\Qbar+\Qbar u+\Qbar\bar u+\Qbar w$.

{\small\emph{Source:} The configurations behind \cite[Corollaire 5(1) and 5(4)]{Dia07}, used there with Roy's strong six exponentials theorem \cite{Roy92}.}
\end{proposition}

\begin{proof}
\leanok
Here $u^{-1}=\bar u/(u\bar u)$ lies in the space. Take $(1,u)\otimes(u^{-1},1,u)$ for $u^{2}$, $(1,u^{-1})\otimes(u,1,u^{-1})$ for $\bar u^{2}$, and $(u^{-1},(u-a)^{-1})\otimes(1,u,u^{2}-au)$ for $1/(u-a)$. The families are free because $u$ is transcendental; the last one uses $a\neq0$.
\end{proof}

% Prove2Me: DiazModulus.circle_point_extension_two_by_three_configuration_iff -- https://prove2.me/theorems/f2c817ff-c9ef-49ae-94c0-78ef031172be
\begin{theorem}[The four-dimensional extensions that carry a configuration]
\label{thm:circle_point_extension_two_by_three_configuration_iff}
\lean{Diaz.circle_point_extension_two_by_three_configuration_iff}
\leanok
Let $u\in\C\setminus\Qbar$ with $u\bar u$ algebraic, put $H_0=\Qbar+\Qbar u+\Qbar\bar u$, and let $z\notin H_0$. There are $x_1,x_2$, linearly independent over $\Qbar$, and $y_1,y_2,y_3$, linearly independent over $\Qbar$, with all six products $x_iy_j$ in $H_0+\Qbar z$ if and only if $z\in H_0+\Qbar w$ for $w=u^{2}$, $w=\bar u^{2}$, or $w=1/(u-a)$ with $a\in\Qbar$, $a\neq0$.

{\small\emph{Source:} \textbf{Not found in the sources read, and not routine} (Chapter~\ref{chap:14-not-found-in-the-sources}). With Roy's strong six exponentials theorem \cite{Roy92}, the ``if'' direction gives at a candidate the exclusions of \cite[Corollaire 5(1) and 5(4)]{Dia07}; every configuration has the shape of \cite[Lemma 6.1]{Fis01} and \cite[Théorème 7(2)]{Dia07}.}
\end{theorem}

\begin{proof}
\leanok
\uses{lem:two_by_three_configuration_forces_progression,lem:circle_hull_progression_contains_square_or_reciprocal,prop:circle_point_extension_carries_two_by_three_configuration}
If $z\in H_0+\Qbar w$ and $z\notin H_0$, then $H_0+\Qbar z=H_0+\Qbar w$, which carries a configuration by Proposition~\ref{prop:circle_point_extension_carries_two_by_three_configuration}. Conversely, Lemma~\ref{lem:two_by_three_configuration_forces_progression} writes $H_0+\Qbar z$ as $b(\Qbar+\Qbar h+\Qbar h^{2}+\Qbar h^{3})$ with $h\notin\Qbar$, and Lemma~\ref{lem:circle_hull_progression_contains_square_or_reciprocal} puts $u^{2}$, $\bar u^{2}$ or some $1/(u-a)$ in it. That element is not in $H_0$, since $u$ is transcendental, so $z\in H_0+\Qbar w$.

Since $\widetilde{\mathcal L}$, the $\Qbar$-span of $1$ and the logarithms of algebraic numbers, is closed under complex conjugation, a hypothesis $z\in\widetilde{\mathcal L}$ also brings $\bar z$; the five-dimensional spaces $H_0+\Qbar z+\Qbar\bar z$ are the subject of the next section.
\end{proof}

\section{Five-dimensional extensions}

% Prove2Me: DiazModulus.circle_point_conjugate_pair_extension_carries_two_by_three_configuration -- https://prove2.me/theorems/d8abac5b-9bd4-4091-a0ed-2be61713f899
\begin{proposition}[Two families of five-dimensional spaces that carry a configuration]
\label{prop:circle_point_conjugate_pair_extension_carries_two_by_three_configuration}
\lean{Diaz.circle_point_conjugate_pair_extension_carries_two_by_three_configuration}
\leanok
Let $u\in\C\setminus\Qbar$ with $\rho=u\bar u$ algebraic, let $a\in\Qbar$, $a\neq0$, and let $z=u/(u^{2}-a)$ if $a\bar a\neq\rho^{2}$, or $z=u/(u^{2}-a)^{2}$ if $a\bar a=\rho^{2}$. Then there are $x_1,x_2$, linearly independent over $\Qbar$, and $y_1,y_2,y_3$, linearly independent over $\Qbar$, with all six products $x_iy_j$ in $\Qbar+\Qbar u+\Qbar\bar u+\Qbar z+\Qbar\bar z$.

{\small\emph{Source:} Rescaled, the configuration is the one in the proofs of \cite[Théorème 2]{Dia04} and \cite[Théorèmes 6(1) and 7(1)]{Dia07}, and the progression of \cite[Lemma 6.1]{Fis01}.}
\end{proposition}

\begin{proof}
\leanok
In the first case put $a'=\rho^{2}/\bar a$, which differs from $a$; then $\bar z=-(\rho/\bar a)\,u/(u^{2}-a')$, and let $b=1/(u(u^{2}-a)(u^{2}-a'))$. In the second, $\bar z=(\rho/\bar a^{2})\,u^{3}/(u^{2}-a)^{2}$, and let $b=1/(u(u^{2}-a)^{2})$. Partial fractions in $u^{2}$, with $u^{-1}=\bar u/\rho$, put $b,bu^{2},bu^{4},bu^{6}$ in the space, so $(1,u^{2})\otimes(b,bu^{2},bu^{4})$ is a configuration, free because $u$ is transcendental.
\end{proof}

% Prove2Me: DiazModulus.circle_point_conjugate_pair_extension_excludes_squares_and_reciprocals -- https://prove2.me/theorems/f640cf63-2fec-48d4-b85e-1eff3f39da34
\begin{lemma}[What these spaces do not contain]
\label{lem:circle_point_conjugate_pair_extension_excludes_squares_and_reciprocals}
\lean{Diaz.circle_point_conjugate_pair_extension_excludes_squares_and_reciprocals}
\leanok
With $u$, $a$ and $z$ as in Proposition~\ref{prop:circle_point_conjugate_pair_extension_carries_two_by_three_configuration}, the space $\Qbar+\Qbar u+\Qbar\bar u+\Qbar z+\Qbar\bar z$ contains neither $u^{2}$ nor $\bar u^{2}$, nor $1/(u-b)$ for any $b\in\Qbar$, $b\neq0$.

{\small\emph{Source:} Elementary.}
\end{lemma}

\begin{proof}
\leanok
Clearing denominators turns a membership into $A(u^{2})+u\,B(u^{2})=0$ with $A,B\in\Qbar[X]$. The same identity holds at $-u$, so $A=B=0$. For $u^{2}$ the odd part is not zero; for $\bar u^{2}=\rho^{2}/u^{2}$ the even part is not zero at $0$; for $1/(u-b)=(u+b)/(u^{2}-b^{2})$ the odd part is not zero at $b^{2}$.
\end{proof}

% Prove2Me: DiazModulus.circle_point_conjugate_pair_configuration_invisible_to_four_dimensional_extensions -- https://prove2.me/theorems/901396cf-2688-4738-be1c-a80b6ff6e7c7
\begin{theorem}[Configurations invisible to the four-dimensional extensions]
\label{thm:circle_point_conjugate_pair_configuration_invisible_to_four_dimensional_extensions}
\lean{Diaz.circle_point_conjugate_pair_configuration_invisible_to_four_dimensional_extensions}
\leanok
With $u$, $a$ and $z$ as in Proposition~\ref{prop:circle_point_conjugate_pair_extension_carries_two_by_three_configuration}, put $H_0=\Qbar+\Qbar u+\Qbar\bar u$ and $W=H_0+\Qbar z+\Qbar\bar z$. Then $W$ carries a configuration, but for no $w\in W\setminus H_0$ does $H_0+\Qbar w$ carry one.

{\small\emph{Source:} Not found in the sources read, but short (Chapter~\ref{chap:14-not-found-in-the-sources}). The shape of the configuration is printed; see Proposition~\ref{prop:circle_point_conjugate_pair_extension_carries_two_by_three_configuration}.}
\end{theorem}

\begin{proof}
\leanok
\uses{prop:circle_point_conjugate_pair_extension_carries_two_by_three_configuration,lem:circle_point_conjugate_pair_extension_excludes_squares_and_reciprocals,thm:circle_point_extension_two_by_three_configuration_iff}
The configuration is Proposition~\ref{prop:circle_point_conjugate_pair_extension_carries_two_by_three_configuration}. If $H_0+\Qbar w$ carried one, Theorem~\ref{thm:circle_point_extension_two_by_three_configuration_iff} would give $w\in H_0+\Qbar t$ with $t$ one of $u^{2}$, $\bar u^{2}$, $1/(u-b)$. Since $w\notin H_0$, exchange puts $t$ in $H_0+\Qbar w\subseteq W$, against Lemma~\ref{lem:circle_point_conjugate_pair_extension_excludes_squares_and_reciprocals}. The configuration does not need all of $W$: its six products span the odd part $\Qbar u+\Qbar\bar u+\Qbar z+\Qbar\bar z$.
\end{proof}

% Prove2Me: DiazModulus.conj_stable_circle_point_extension_no_two_by_three_configuration -- https://prove2.me/theorems/9819e926-24a6-4c12-8431-7a5f6be68dfa
\begin{proposition}[The conjugation-stable case carries none]
\label{prop:conj_stable_circle_point_extension_no_two_by_three_configuration}
\lean{Diaz.conj_stable_circle_point_extension_no_two_by_three_configuration}
\leanok
Let $u\in\C\setminus\Qbar$ with $\rho=u\bar u$ algebraic, and let $a\in\Qbar$, $a\neq0$, with $a\bar a=\rho^{2}$. Then there are no $x_1,x_2$, linearly independent over $\Qbar$, and $y_1,y_2,y_3$, linearly independent over $\Qbar$, with all six products $x_iy_j$ in $\Qbar+\Qbar u+\Qbar\bar u+\Qbar\,u/(u^{2}-a)$.

{\small\emph{Source:} A consequence of Theorem~\ref{thm:circle_point_conjugate_pair_configuration_invisible_to_four_dimensional_extensions}. It is where the hypotheses of \cite[Théorème 2]{Dia04} and of \cite[Corollaire 4(4) and Théorème 7(1)]{Dia07} fail for this family.}
\end{proposition}

\begin{proof}
\leanok
\uses{thm:circle_point_conjugate_pair_configuration_invisible_to_four_dimensional_extensions}
Here $z=u/(u^{2}-a)$ has $\bar z=-(\rho/\bar a)z$, so the space is stable under conjugation. With $z_2=u/(u^{2}-a)^{2}$, the element $u/(u^{2}-a)=(\bar a^{2}/\rho)\bar z_2-az_2$ lies in $H_0+\Qbar z_2+\Qbar\bar z_2$ and not in $H_0$, so Theorem~\ref{thm:circle_point_conjugate_pair_configuration_invisible_to_four_dimensional_extensions} applies. When $a\bar a\neq\rho^{2}$, adding $\bar z$ gives a space that does carry a configuration.
\end{proof}

% Prove2Me: DiazModulus.logAlg_conj_stable -- https://prove2.me/theorems/95246d4a-a61c-4120-bfd0-0f30ffec3055
\begin{lemma}[Conjugation keeps the logarithms]
\label{lem:logAlg_conj_stable}
\lean{Diaz.logAlg_conj_stable}
\leanok
If $\lambda$ is a logarithm of an algebraic number, so is $\bar\lambda$.

{\small\emph{Source:} Elementary.}
\end{lemma}

\begin{proof}
\leanok
$e^{\bar\lambda}=\overline{e^{\lambda}}$, and the conjugate of an algebraic number is algebraic.
\end{proof}

% Prove2Me: DiazModulus.logAlgTilde_conj_stable -- https://prove2.me/theorems/9922f109-5be0-427a-93a2-e7015b8cdf9d
\begin{lemma}[Conjugation keeps $\widetilde{\mathcal L}$]
\label{lem:logAlgTilde_conj_stable}
\lean{Diaz.logAlgTilde_conj_stable}
\leanok
The $\Qbar$-vector space $\widetilde{\mathcal L}$ spanned by $1$ and the logarithms of algebraic numbers is closed under complex conjugation.

{\small\emph{Source:} Elementary; it is the stability that \cite[Théorème 2]{Dia04} uses.}
\end{lemma}

\begin{proof}
\leanok
\uses{lem:logAlg_conj_stable}
Conjugation fixes $1$, sends logarithms to logarithms (Lemma~\ref{lem:logAlg_conj_stable}), and $\overline{c\lambda}=\bar c\,\bar\lambda$ with $\bar c\in\Qbar$.
\end{proof}

% Prove2Me: DiazModulus.candidate_div_sq_sub_not_mem_logAlgTilde -- https://prove2.me/theorems/35486890-e3f7-4a35-a512-441fc709f738
\begin{corollary}[Two exclusions at a candidate]
\label{cor:candidate_div_sq_sub_not_mem_logAlgTilde}
\lean{Diaz.candidate_div_sq_sub_not_mem_logAlgTilde}
\leanok
Assume Roy's strong six exponentials theorem, and let $u$ be a candidate: $u\neq0$, $|u|$ algebraic and $e^{u}$ algebraic. Let $a\in\Qbar$, $a\neq0$, and $\rho=u\bar u$. Then $u/(u^{2}-a)\notin\widetilde{\mathcal L}$ if $a\bar a\neq\rho^{2}$, and $u/(u^{2}-a)^{2}\notin\widetilde{\mathcal L}$ if $a\bar a=\rho^{2}$.

{\small\emph{Source:} One substitution in \cite[Théorème 2]{Dia04}, at $x=(u,\bar u)$ and $y=1/(u^{2}-a)$, and in \cite[Corollaire 4(4) and Théorème 7(1)]{Dia07}; the second part is \cite[Théorème 6(3)]{Dia07}. Roy's theorem is \cite{Roy92}.}
\end{corollary}

\begin{proof}
\leanok
\uses{prop:circle_point_conjugate_pair_extension_carries_two_by_three_configuration,thm:hermite_lindemann_holds,lem:logAlg_conj_stable,lem:logAlgTilde_conj_stable}
$\widetilde{\mathcal L}$ is closed under conjugation and contains $1$, $u$ and $\bar u$, so $z\in\widetilde{\mathcal L}$ would put the space of Proposition~\ref{prop:circle_point_conjugate_pair_extension_carries_two_by_three_configuration} in $\widetilde{\mathcal L}$, with its configuration, against Roy's theorem.
\end{proof}

\section{Power hulls}

% Prove2Me: DiazModulus.power_hull_strong_six_exp_configuration_iff -- https://prove2.me/theorems/e31ad15c-0c6d-4572-b443-953da5c153e1
\begin{theorem}[Configurations in the power hulls]
\label{thm:power_hull_strong_six_exp_configuration_iff}
\lean{Diaz.power_hull_strong_six_exp_configuration_iff}
\leanok
Let $u\in\C\setminus\Qbar$ and $k\ge1$. There are $x_1,x_2$, linearly independent over $\Qbar$, and $y_1,y_2,y_3$, linearly independent over $\Qbar$, with all six products $x_iy_j$ in $\Qbar u^{-k}+\Qbar u^{-1}+\Qbar+\Qbar u+\Qbar u^{k}$ if and only if $k=2$ or $k=3$.

{\small\emph{Source:} \textbf{Not found in the sources read, and not routine} (Chapter~\ref{chap:14-not-found-in-the-sources}). For $k=2,3$ the configurations are geometric progressions of four elements, which \cite[Lemma 6.1]{Fis01} and \cite[Théorème 7(2)]{Dia07} exclude from $\widetilde{\mathcal L}$ with Roy's theorem; at a candidate they give \cite[Corollaire 5(1) and 5(2)]{Dia07}. The direction ``only for $k=2,3$'' was not found; Diaz remarks that inside that theorem one cannot hope to go very far \cite[p.~390]{Dia07}.}
\end{theorem}

\begin{proof}
\leanok
For $k=2$ take $(1,u)\otimes(u^{-1},1,u)$, and for $k=3$ take $(u,u^{-1})\otimes(1,u^{-2},u^{2})$. Conversely, multiply by $u^{k}$: the products become polynomials in $u$ supported on $\{0,k-1,k,k+1,2k\}$, and since $u$ is transcendental a rank-one $2\times3$ matrix of them, with free rows and columns, has rows that differ by a fixed non-zero shift of the orders at $0$. For $k=1$ and for $k\ge4$ no shift keeps three exponents of that set inside it.
\end{proof}
